% AUTO-GENERATED. Edit the Markdown sources or notes.config.json.
% Sources: 02-Economics/2026-Financial-Economics/2026-09-24-Note.md

\section{货币的时间价值}

\subsection{The Time Value of Money}

\textbf{Definition}

\begin{quote}
The \textbf{time value of money} means that cash flows received at different dates cannot be compared or added directly without first moving them to the same date.
\end{quote}

今天的 1 元通常比未来确定收到的 1 元更有价值，原因包括：

\begin{itemize}
\tightlist
\item
  当前资金可以投资并产生利息；
\item
  通货膨胀会降低未来货币的购买力；
\item
  未来现金流可能存在违约或其他不确定性。
\end{itemize}

金融计算的基本原则是：先选定一个 valuation date，再通过 compounding 或 discounting 把所有现金流移动到该日期。

\begin{longtable}[]{@{}>{\raggedright\arraybackslash}p{0.22\linewidth}>{\raggedright\arraybackslash}p{0.34\linewidth}>{\raggedright\arraybackslash}p{0.34\linewidth}@{}}
\toprule\noalign{}
符号 & English term & 中文含义 \\
\midrule\noalign{}
\endhead
\bottomrule\noalign{}
\endlastfoot
\(PV\) & present value & 现值 \\
\(FV\) & future value & 终值 \\
\(i\) & periodic interest rate & 每期利率 \\
\(n\) & number of periods & 计息期数 \\
\(PMT\) & periodic payment & 每期等额现金流 \\
\end{longtable}

\subsection{Lump-Sum Cash Flows}

\subsubsection{Future value}

一笔现值为 \(PV\) 的资金按每期利率 \(i\) 复利增长 \(n\) 期，终值为

\[
\boxed{FV=PV(1+i)^n}.
\]

例如，将 \(1{,}500\) 元按年利率 \(3\%\) 投资 5 年：

\[
FV=1500(1.03)^5\approx1{,}738.91.
\]

复利的关键在于每一期的利息也会进入下一期本金。

\subsubsection{Present value}

把未来现金流贴现到今天：

\[
\boxed{
PV=\frac{FV}{(1+i)^n}=FV(1+i)^{-n}
}.
\]

若两年后收到 \(40{,}000\) 元，要求回报率为 \(8\%\)，则

\[
PV=\frac{40{,}000}{(1.08)^2}
\approx34{,}293.55.
\]

要求回报率越高，给定未来现金流的现值越低。

\subsubsection{Solving for the interest rate}

由

\[
FV=PV(1+i)^n
\]

可得

\[
\boxed{
i=\left(\frac{FV}{PV}\right)^{1/n}-1
}.
\]

例如，\(15{,}000\) 元在十年后增长为 \(30{,}000\) 元，则年化复合收益率为

\[
i=2^{1/10}-1\approx7.18\%.
\]

\subsubsection{Solving for the number of periods}

对终值公式取自然对数：

\[
\ln\frac{FV}{PV}=n\ln(1+i),
\]

因此

\[
\boxed{
n=\frac{\ln(FV/PV)}{\ln(1+i)}
}.
\]

\subsubsection{Rule of 72}

资金翻倍所需年数可近似为

\[
\text{doubling time}\approx\frac{72}{\text{annual rate in percent}}.
\]

这一规则适合快速估算，不能替代精确的复利计算。精确翻倍时间为

\[
n=\frac{\ln2}{\ln(1+i)}.
\]

\begin{quote}
\textbf{注：Calculation precision}

中间步骤应保留足够精度，只在最终结果处按经济意义进行舍入。过早截断会在长期复利和大额项目中积累明显误差。
\end{quote}

\subsection{Compounding Frequency}

设名义年利率，即 annual percentage rate 为 \(k\)，每年复利 \(m\) 次，则每期利率为

\[
\frac{k}{m}.
\]

投资 1 元一年后的价值为

\[
\left(1+\frac{k}{m}\right)^m.
\]

因此有效年利率为

\[
\boxed{
EFF=\left(1+\frac{k}{m}\right)^m-1
}.
\]

\subsubsection{APR 与有效年利率}

\textbf{Definition}

\begin{quote}
The \textbf{annual percentage rate, APR}, reports a nominal annual rate without incorporating within-year compounding into the quoted percentage.
\end{quote}

\textbf{Definition}

\begin{quote}
The \textbf{effective annual rate, EFF}, is the one-year return after accounting for all compounding within the year.
\end{quote}

例如，信用卡 APR 为 \(18\%\)，按月复利，则月利率为 \(1.5\%\)，有效年利率为

\[
EFF=(1.015)^{12}-1\approx19.56\%.
\]

相同 APR 下，复利频率越高，EFF 越高。因此比较不同贷款或存款报价时，应使用相同期间的有效利率。

\subsubsection{从有效年利率反推 APR}

若目标有效年利率为 \(EFF\)，每年复利 \(m\) 次，则对应 APR 为

\[
\boxed{
k_m=m\left[(1+EFF)^{1/m}-1\right]
}.
\]

例如，银行希望获得 \(12\%\) 的有效年利率时，复利越频繁，对应的报价 APR 越低，但实际一年增长因子保持 \(1.12\) 不变。

\subsubsection{Continuous compounding}

当复利频率趋于无穷，

\[
\lim_{m\to\infty}
\left(1+\frac{k}{m}\right)^m=e^k.
\]

连续复利下的有效年利率为

\[
\boxed{EFF=e^k-1}.
\]

APR 为 \(18\%\) 且连续复利时，

\[
EFF=e^{0.18}-1\approx19.72\%.
\]

\subsection{Cash-Flow Timelines}

时间线将每笔现金流放在其实际发生日期。若 \(CF_t\) 表示第 \(t\) 期末现金流，则其现值为

\[
\boxed{
PV=\sum_{t=0}^{n}\frac{CF_t}{(1+i)^t}
}.
\]

在第 \(n\) 期计算终值时，

\[
\boxed{
FV_n=\sum_{t=0}^{n}CF_t(1+i)^{n-t}
}.
\]

\textbf{Value additivity} 表明，同一估值日期上的现金流价值可以相加：

\[
PV(CF^{(1)}+CF^{(2)})
=PV(CF^{(1)})+PV(CF^{(2)}).
\]

不同日期的名义金额不能直接相加，必须先移动到同一日期。

例如，若贴现率为 \(10\%\)，第 1 期收到 \(1{,}000\) 元、第 2 期收到 \(2{,}000\) 元，则

\[
PV=\frac{1{,}000}{1.10}+\frac{2{,}000}{(1.10)^2}
\approx909.09+1{,}652.89
=2{,}561.98.
\]

\subsubsection{Net Present Value and Internal Rate of Return}

\textbf{Definition}

\begin{quote}
The \textbf{net present value, NPV}, is the present value of all project cash inflows minus the present value of all project cash outflows.
\end{quote}

若 \(CF_0\) 已包含初始投资，则

\[
\boxed{
NPV(r)=\sum_{t=0}^{n}\frac{CF_t}{(1+r)^t}
}.
\]

在贴现率已经正确反映资金机会成本和项目风险时，独立项目的决策规则是：

\[
NPV>0\Rightarrow\text{accept},
\qquad
NPV<0\Rightarrow\text{reject}.
\]

\textbf{Definition}

\begin{quote}
The \textbf{internal rate of return, IRR}, is a discount rate that makes the project's NPV equal to zero.
\end{quote}

即 IRR 满足

\[
0=\sum_{t=0}^{n}\frac{CF_t}{(1+IRR)^t}.
\]

对先发生一次净流出、随后只产生净流入的 conventional cash flows，若 \(IRR\) 高于 required rate of return，项目通常可以接受。但存在多次现金流符号变化时，可能出现多个 IRR 或不存在实数 IRR；互斥项目的规模和时点不同也可能使 IRR 排序与 NPV 排序冲突。因此，最终价值判断应以 NPV 为主。

\subsection{Ordinary Annuities}

\textbf{Definition}

\begin{quote}
An \textbf{ordinary annuity} is a sequence of equal cash flows occurring at equally spaced dates, with the first payment one period from today.
\end{quote}

标准年金公式通常假设：

\begin{itemize}
\tightlist
\item
  每期现金流均为 \(PMT\)；
\item
  每期长度相同；
\item
  每期利率恒定为 \(i\)；
\item
  第一次支付发生在一期后。
\end{itemize}

\subsubsection{Present value of an annuity}

普通年金现值为

\[
PV_a
=\frac{PMT}{1+i}
+\frac{PMT}{(1+i)^2}
+\cdots
+\frac{PMT}{(1+i)^n}.
\]

利用等比数列可得

\[
\boxed{
PV_a=PMT\frac{1-(1+i)^{-n}}{i}
}.
\]

年金现值因子为

\[
PVIFA(i,n)=\frac{1-(1+i)^{-n}}{i}.
\]

\subsubsection{Future value of an annuity}

普通年金在第 \(n\) 期末的终值为

\[
\boxed{
FV_a=PMT\frac{(1+i)^n-1}{i}
}.
\]

\subsubsection{Solving for the payment}

给定贷款或年金现值时，每期支付额为

\[
\boxed{
PMT=PV\frac{i}{1-(1+i)^{-n}}
}.
\]

\subsubsection{Solving for the number of payments}

由普通年金现值公式可得

\[
\boxed{
n=-\frac{
\ln\left(1-\dfrac{PV\cdot i}{PMT}\right)
}{\ln(1+i)}
},
\]

前提是 \(PMT>PV\cdot i\)。如果每期支付不超过当期利息，本金无法在有限期内还清。

给定 \(PV\)、\(PMT\) 和 \(n\) 反求利率 \(i\) 时，一般不存在简单闭式解，需要使用数值方法或 financial calculator。

\subsection{Annuities Due}

\textbf{Definition}

\begin{quote}
An \textbf{annuity due} makes its first payment today rather than one period from today.
\end{quote}

Annuity due 的每笔现金流都比普通年金提前一期，因此

\[
\boxed{
PV_{ad}=PV_a(1+i)
},
\]

\[
\boxed{
FV_{ad}=FV_a(1+i)
}.
\]

判断年金类型时，不要只看支付次数，要先确定第一笔支付发生在 \(t=0\) 还是 \(t=1\)。

\begin{quote}
\textbf{警告：术语核对}

课件第 41 页把从 \(t=0\) 开始支付的年金标为 ``Immediate Annuity''。标准金融与精算术语通常将其称为 \textbf{annuity due}；\textbf{annuity-immediate} 通常指第一笔付款发生在一期后的普通年金。
\end{quote}

\subsection{Perpetuities}

\subsubsection{Level perpetuity}

若每期永久支付固定金额 \(PMT\)，且 \(i>0\)，则

\[
\boxed{
PV=\frac{PMT}{i}
}.
\]

这是普通年金现值在 \(n\to\infty\) 时的极限。

\subsubsection{Growing perpetuity}

若下一期的第一笔现金流为 \(PMT_1\)，以后每期按固定速度 \(g\) 永久增长，且 \(i>g\)，则

\[
\boxed{
PV=\frac{PMT_1}{i-g}
}.
\]

\begin{quote}
\textbf{警告：术语核对}

课件第 56 页将 \(PV=PMT_1/(i-g)\) 放在 ``Growing Annuities'' 标题下。该公式对应 \textbf{growing perpetuity}；有限期 growing annuity 还需要包含终止期 \(n\)。
\end{quote}

\subsection{Loan Amortization}

\textbf{Definition}

\begin{quote}
\textbf{Loan amortization} is the gradual repayment of loan principal through a sequence of periodic payments containing both interest and principal.
\end{quote}

固定利率、等额支付贷款的每期还款额为

\[
PMT=PV\frac{i}{1-(1+i)^{-n}}.
\]

每一期均满足

\[
\begin{aligned}
\text{interest}_t&=i\times\text{balance}_{t-1},\\
\text{principal}_t&=PMT-\text{interest}_t,\\
\text{balance}_t&=\text{balance}_{t-1}-\text{principal}_t.
\end{aligned}
\]

随着未偿本金下降，利息部分逐期减少，本金部分逐期增加。

\subsubsection{Mortgage example}

房价为 \(1{,}000{,}000\) 元，首付比例为 \(10\%\)，贷款期限为 360 个月，月利率为 \(0.5\%\)：

\[
PV=1{,}000{,}000(1-0.10)=900{,}000.
\]

Three discount points 等于贷款额的 \(3\%\)：

\[
\text{points}=900{,}000\times0.03=27{,}000.
\]

借款人初始实际获得的净资金为

\[
900{,}000-27{,}000=873{,}000,
\]

但仍按 \(900{,}000\) 的名义本金偿还贷款，因此 points 会提高贷款的有效融资成本。

每月还款额为

\[
\begin{aligned}
PMT
&=900{,}000\frac{0.005}{1-(1.005)^{-360}}\\
&\approx5{,}395.95.
\end{aligned}
\]

第一个月：

\[
\begin{aligned}
\text{interest}_1&=900{,}000\times0.005=4{,}500,\\
\text{principal}_1&=5{,}395.95-4{,}500=895.95,\\
\text{balance}_1&=900{,}000-895.95=899{,}104.05.
\end{aligned}
\]

\begin{quote}
\textbf{警告：课件数值核对}

第 62 页把首月末余额写成 \(889{,}104.04\)，正确结果约为 \(899{,}104.05\)。这是减法中的位数错误。
\end{quote}

\subsubsection{Outstanding balance after 60 payments}

支付 60 期后，还剩 300 期。第 60 期末的未偿本金等于剩余 300 笔付款在该时点的现值：

\[
\begin{aligned}
B_{60}
&=PMT\frac{1-(1.005)^{-300}}{0.005}\\
&\approx837{,}489.21.
\end{aligned}
\]

前 60 期总付款约为

\[
5{,}395.9547\times60\approx323{,}757.28.
\]

其中本金减少

\[
900{,}000-837{,}489.21=62{,}510.79,
\]

利息支出约为

\[
323{,}757.28-62{,}510.79=261{,}246.49.
\]

前期付款中利息占比较高，是因为利息按期初未偿余额计算，而初期本金仍然很大。

\subsubsection{Early repayment decisions}

判断是否提前还贷时，不能把未来几十年的名义利息简单相加并称为``节省''。正确比较需要：

\begin{enumerate}
\def\labelenumi{\arabic{enumi}.}
\tightlist
\item
  把提前还款与继续还款的现金流放在时间线上；
\item
  使用与风险相匹配的贴现率计算现值；
\item
  考虑税收、提前还款费用和替代投资机会；
\item
  比较风险调整后的净现值，而非未贴现现金总额。
\end{enumerate}

\subsection{Exchange Rates and Time Value of Money}

跨币种估值必须满足：现金流和贴现率使用同一种货币。可以先在各自货币下计算现值，再用估值日即期汇率换算。

\subsubsection{Two projects}

美国项目 \(U\) 初始投资 \(10{,}000\) 美元，未来 5 年每年产生 \(6{,}000\) 美元，美元利率为 \(6\%\)：

\[
\begin{aligned}
NPV_U
&=6{,}000\frac{1-(1.06)^{-5}}{0.06}-10{,}000\\
&\approx15{,}274.18\ \text{USD}.
\end{aligned}
\]

日本项目 \(J\) 初始投资 \(1{,}000{,}000\) 日元，未来 5 年每年产生 \(575{,}000\) 日元，日元利率为 \(4\%\)：

\[
\begin{aligned}
NPV_J
&=575{,}000\frac{1-(1.04)^{-5}}{0.04}-1{,}000{,}000\\
&\approx1{,}559{,}797.84\ \text{JPY}.
\end{aligned}
\]

若即期汇率为 \(0.01\) USD/JPY，则

\[
NPV_J\approx15{,}597.98\ \text{USD}.
\]

在忽略汇率风险和其他差异的条件下，项目 \(J\) 的现值略高。但如果未来现金流最终要换成美元，仍需分析汇率风险或使用远期合约套期保值。

\begin{quote}
\textbf{警告：课件数值核对}

第 74--75 页的日元 NPV 排版存在多余数字。按课件参数计算，应约为 \(1{,}559{,}797.84\) 日元，即 \(15{,}597.98\) 美元。
\end{quote}

\subsection{Inflation and Discounted Cash Flow}

设名义利率为 \(i_n\)，实际利率为 \(i_r\)，通货膨胀率为 \(\pi\)。精确 Fisher relation 为

\[
\boxed{
1+i_n=(1+i_r)(1+\pi)
}.
\]

因此

\[
\boxed{
i_r=\frac{i_n-\pi}{1+\pi}
}.
\]

低通胀下常用近似为

\[
i_r\approx i_n-\pi.
\]

若名义利率为 \(8\%\)、通胀率为 \(5\%\)，则

\[
i_r=\frac{0.08-0.05}{1.05}\approx2.86\%.
\]

\subsubsection{Nominal consistency rule}

估值时必须保持口径一致：

\[
\text{nominal cash flows}
\longleftrightarrow
\text{nominal discount rate},
\]

\[
\text{real cash flows}
\longleftrightarrow
\text{real discount rate}.
\]

混用实际现金流与名义贴现率，或名义现金流与实际贴现率，会系统性扭曲现值。

\subsubsection{Nominal security versus inflation-indexed security}

若名义存款收益率为 \(8\%\)，另一项资产提供 \(3\%\) 的实际收益率，则精确的 break-even inflation 满足

\[
1.08=1.03(1+\pi^*).
\]

因此

\[
\pi^*=\frac{1.08}{1.03}-1\approx4.85\%.
\]

用 \(i_n\approx i_r+\pi\) 的近似计算时，临界通胀约为 \(5\%\)。

\subsubsection{Unexpected inflation and borrowers}

固定名义利率贷款的实际偿还成本为

\[
\frac{1+i_n}{1+\pi}-1.
\]

实际通胀高于签约时预期，会降低债务的实际价值，因此有利于借款人、不利于贷款人；实际通胀低于预期则相反。

\subsection{Taxes and After-Tax Returns}

若利息收入按边际税率 \(\tau\) 纳税，且没有其他抵扣，税后利率为

\[
\boxed{
i_{\text{after-tax}}=i_{\text{before-tax}}(1-\tau)
}.
\]

投资决策应最大化风险调整后的税后收益，而不是单纯最小化纳税额。

如果免税债券收益率为 \(i_m\)，应税债券的税前收益率至少要达到

\[
\boxed{
i_{\text{taxable equivalent}}
=\frac{i_m}{1-\tau}
}.
\]

例如，免税债券收益率为 \(6\%\)，边际税率为 \(20\%\)，则无差异的应税收益率为

\[
\frac{0.06}{1-0.20}=7.5\%.
\]

\subsection{复习要点}

\begin{enumerate}
\def\labelenumi{\arabic{enumi}.}
\tightlist
\item
  不同日期的现金流必须先移动到同一估值日期，才能比较或相加。
\item
  APR 不包含年内复利效果；比较报价时应换算为 EFF。
\item
  普通年金第一笔付款在一期后，annuity due 第一笔付款在今天，后者价值是前者的 \((1+i)\) 倍。
\item
  永续年金公式要求 \(i>0\)；增长永续年金公式要求 \(i>g\)。
\item
  贷款还款额由利息与本金构成，未偿本金等于剩余付款在当前时点的现值。
\item
  提前还贷不能依据未贴现利息总额判断，必须比较两套现金流的现值与机会成本。
\item
  跨币种估值时，现金流与贴现率必须使用相同货币。
\item
  名义现金流对应名义贴现率，实际现金流对应实际贴现率。
\item
  税收比较应使用税后收益或 taxable-equivalent yield。
\end{enumerate}

\subsection{术语表}

\begin{longtable}[]{@{}
  >{\raggedright\arraybackslash}p{(\linewidth - 4\tabcolsep) * \real{0.3333}}
  >{\raggedright\arraybackslash}p{(\linewidth - 4\tabcolsep) * \real{0.3333}}
  >{\raggedright\arraybackslash}p{(\linewidth - 4\tabcolsep) * \real{0.3333}}@{}}
\toprule\noalign{}
\begin{minipage}[b]{\linewidth}\raggedright
Term
\end{minipage} & \begin{minipage}[b]{\linewidth}\raggedright
中文
\end{minipage} & \begin{minipage}[b]{\linewidth}\raggedright
简要含义
\end{minipage} \\
\midrule\noalign{}
\endhead
\bottomrule\noalign{}
\endlastfoot
Compounding & 复利 & 将利息计入本金并继续产生利息 \\
Discounting & 贴现 & 将未来现金流换算为当前价值 \\
Annual percentage rate & 名义年利率 & 未把年内复利效果计入报价百分比的年利率 \\
Effective annual rate & 有效年利率 & 计入一年内全部复利后的实际年增长率 \\
Ordinary annuity & 普通年金 & 第一次付款发生在一期后的等额现金流序列 \\
Annuity due & 期初年金 & 第一次付款发生在今天的等额现金流序列 \\
Perpetuity & 永续年金 & 无限期持续支付的现金流序列 \\
Growing perpetuity & 增长永续年金 & 按固定增长率永久增长的现金流序列 \\
Loan amortization & 贷款摊销 & 通过定期付款逐步偿还贷款本金 \\
Discount points & 贷款点数 & 按贷款本金比例在放款时收取的预付费用 \\
Fisher relation & Fisher 关系 & 名义利率、实际利率与通胀率之间的精确关系 \\
Taxable-equivalent yield & 应税等价收益率 & 与免税收益率产生相同税后回报的应税收益率 \\
\end{longtable}
